\pdfminorversion=7%
\documentclass[aspectratio=169,mathserif,notheorems]{beamer}%
%
\xdef\bookbaseDir{../../bookbase}%
\xdef\sharedDir{../../shared}%
\RequirePackage{\bookbaseDir/styles/slides}%
\RequirePackage{\sharedDir/styles/styles}%
%
\gdef\searchSpace{\ensuremath{\mathbb{X}}}%
%
\protected\gdef\algobox#1#2{\resizebox{#1\paperwidth}{!}{%
\bgroup%
\fboxsep=0pt%
\fboxrule=2pt%
\mbox{\fcolorbox{white}{white}{\bgroup%
\fboxsep=2pt%
\fboxrule=2pt%
\mbox{\fcolorbox{hfuu-red!80}{hfuu-orange!60}{%
\mbox{#2}%
}}\egroup%
}}%
\egroup}}%
%
\title{Metaheuristic Optimization: State of the Art}%
%
%
\begin{document}%
\startPresentation%
%
\section{Introduction}%
%
\begin{frame}%
\frametitle{Metaheuristics}%
\begin{itemize}%
\item Metaheuristics are a class of algorithms used for \alert{optimization}.%
\end{itemize}%
\end{frame}%
%
%
\begin{frame}%
\sharedframetitle{What is Optimization?}%
%
\begin{itemize}%
\item There are \alert{three} ways to approach this topic\only<-1>{.}\uncover<2->{\medskip%
\begin{enumerate}%
%
\item Optimization means searching for \alert{superlatives}.\medskip%
%
\item<3-> Optimization is the art of solving \alert{hard} problems.\medskip%
%
\item<4-> Optimization allows us to find acceptable solutions for problems where our tools from high school maths \emph{and} traditional algorithms from computer science are no longer enough, problems where exact solutions \alert{cannot always be reached}.%
%
\end{enumerate}}%
\end{itemize}%
%
\end{frame}%
%
\global\let\sharedSlideTitle\relax%
%
%
%
\begin{frame}%
\frametitle{Views on Optimization, Part~1}%
\begin{itemize}%
\item Let's start with the first view: \alert{Superlatives}.%
\end{itemize}%
\end{frame}%
%
\input{\sharedDir/slides/optimization_superlatives.tex}%
%
\begin{frame}%
\frametitle{Views on Optimization, Part~2}%
\begin{itemize}%
\item Now let us consider the second view: \alert{Solving hard problems}.%
\item<2-> Later, we learn how to quantify \emph{hard}.%
\item<3-> For now, let's first think about how to define \emph{problem}?%
\end{itemize}%
\end{frame}%
%
%
\begin{frame}[t]%
\sharedframetitle{What are Optimization Problems?}%
\begin{itemize}%
\only<-1>{\item There are two ways to define optimization problems.}%
\only<-2>{\item<2-> The economic view.}%
\item<3-> The mathematical view.%
\end{itemize}%
\locateGraphic{2}{width=0.86\paperwidth}{\sharedDir/graphics/optimization_views/optimization_views_1}{0.07}{0.25}%
\locateGraphic{3}{width=0.86\paperwidth}{\sharedDir/graphics/optimization_views/optimization_views_2}{0.07}{0.25}%
\end{frame}%
%
\global\let\sharedSlideTitle\relax%
%
%
%
\begin{frame}%
\sharedframetitle{Optimization in Business Settings}%
%
\locateGraphic{1}{width=1.05\paperwidth}{\sharedDir/graphics/five_aspects/five_aspects_1}{-0.1}{0.05}%
\locateGraphic{2}{width=1.05\paperwidth}{\sharedDir/graphics/five_aspects/five_aspects_2}{-0.1}{0.05}%
\locateGraphic{3}{width=1.05\paperwidth}{\sharedDir/graphics/five_aspects/five_aspects_3}{-0.1}{0.05}%
\locateGraphic{4}{width=1.05\paperwidth}{\sharedDir/graphics/five_aspects/five_aspects_4}{-0.1}{0.05}%
\locateGraphic{5}{width=1.05\paperwidth}{\sharedDir/graphics/five_aspects/five_aspects_5}{-0.1}{0.05}%
\locateGraphic{6}{width=1.05\paperwidth}{\sharedDir/graphics/five_aspects/five_aspects_6}{-0.1}{0.05}%
\locateGraphic{7}{width=1.05\paperwidth}{\sharedDir/graphics/five_aspects/five_aspects}{-0.1}{0.05}%
%
\end{frame}%
%
%
%
\begin{frame}%
\frametitle{Views on Optimization, Part~3}%
\begin{itemize}%
%
\item This brings us to the third view of optimization: We need to solve problems where equations are \inQuotes{not enough} and for which we cannot always find the best solution.%
%
\item<2-> Let's start with high school Maths, move to undergraduate computer science, and then approach this field of problems.%
\end{itemize}%
\end{frame}%
%
\definecolor{square-color}{HTML}{A03214}%
\definecolor{rectangle-color}{HTML}{0000A0}%
\definecolor{small-triangle-color}{HTML}{AA0078}%
\definecolor{large-triangle-color}{HTML}{008200}%
\definecolor{x-color}{HTML}{00788C}%
%
\protected\gdef\aRectangle{{\color{rectangle-color}{\ensuremath{A_{\text{\ding{122}}}}}}}%
\protected\gdef\aRectangleDiff{{\color{rectangle-color}{\ensuremath{A'_{\text{\ding{122}}}}}}}%
\protected\gdef\aRectangleMax{{\color{rectangle-color}{\ensuremath{\widehat{A_{\text{\ding{122}}}}}}}}%
\protected\gdef\aSquare{{\color{square-color}{\ensuremath{A_{\text{\ding{110}}}}}}}%
\protected\gdef\aSmallTriangle{{\color{small-triangle-color}{\ensuremath{A_{\text{\ding{116}}}}}}}%
\protected\gdef\aLargeTriangle{{\color{large-triangle-color}{\ensuremath{A_{\text{\ding{115}}}}}}}%
\protected\gdef\aX{{\color{x-color}{\ensuremath{x}}}}%
\protected\gdef\aXMax{{\color{x-color}{\ensuremath{\hat{x}}}}}%
\protected\gdef\aRectangleOfX{\ensuremath{\aRectangle(\aX)}}%
\protected\gdef\aRectangleOfXDiff{\ensuremath{\aRectangleDiff(\aX)}}%
%
\begin{frame}[t]%
\frametitle{A Problem that we can solve with high school Maths equations}%
\noindent\parbox{0.65\linewidth}{\noindent%
\begin{itemize}%
\item A \textcolor{rectangle-color}{rectangle} should be placed inside a \textcolor{square-color}{square} with side\\length~13cm as sketched.%
%
\item<2-> What area has the \emph{largest} such \textcolor{rectangle-color}{rectangle}?%
%
\item<3-> We can solve this by realizing that $\aRectangleOfX={\color{square-color}{\left(13\text{cm}^2\right)}}-2{\color{small-triangle-color}{\left(\tfrac{1}{2}\aX^2\right)}}-2{\color{large-triangle-color}{\left(\tfrac{1}{2}(13\text{cm}-\aX)^2\right)}}$\only<-3>{.}\uncover<4->{ and then finding the maximum by setting $0=\aRectangleOfXDiff=-4\aX+26\text{cm}$.}%
%
\item<5-> We get $\aXMax=6.5\text{cm}$\only<-5>{.}\uncover<6->{ and, hence, $\aRectangleMax=84.5\text{cm}^2$.}%
%
\item<7-> Maxima and minima are superlatives $=$ optima.%
%
\item<8-> Sometimes we can simply compute them and do not need fancy algorithms.%
\end{itemize}}%
%%
\strut\\[\paperheight]\strut%
%
\locateGraphic{1-2}{width=0.3\paperwidth}{\sharedDir/graphics/rectangle_problem/rectangle_problem_1}{0.67}{0.15}%
\locateGraphic{3-}{width=0.3\paperwidth}{\sharedDir/graphics/rectangle_problem/rectangle_problem_3}{0.67}{0.15}%
%
\end{frame}%
%

\definecolor{start-color}{HTML}{FF00C8}%
\definecolor{goal-color}{HTML}{00C0C0}%
\definecolor{path-red}{HTML}{FF0000}%
\definecolor{path-pink}{HTML}{FF00FF}%
\definecolor{path-green}{HTML}{008000}%
\definecolor{path-cyan}{HTML}{00C8C8}%
\definecolor{path-gray}{HTML}{646464}%
\protected\gdef\apath#1#2{{\color{path-#1}{#2}}}%
\protected\gdef\noapath{\protected\gdef\apath##1##2{##2}}%
%
\begin{frame}[t]%
\frametitle{Problems that algorithms can solve efficiently}%
%
\locate{2-}{\parbox{0.47\paperwidth}{\noindent%
\begin{itemize}%
%
\only<-9>{%
\item I am at the {\color{start-color}{starting point 合肥大学南二区北大门}}.%
%
\item<3-> I want to go to the {\color{goal-color}{goal:~食堂}}.
%
\item<4-> How do I get there the \alert{fastest}?%
%
\item<5-> We know the campus map.%
%
\item<6-> We want to compute the shortest path~(before actually walking it).%
%
\item<7-> We know all the intersections where I could make turns.%
%
\item<8-> There is no mathematical equation that we can use to solve this.%
}%
%
\item<9-> But we can use the \pgls{algoAstar}\cite{HNR1968AFBFTHDOMCP,P1997AAP}, which iteratively constructs a solution by combining a {\color{path-green}{heuristic}} and a {\color{path-red}{cost}} function.%
%
\only<-15>{%
\item<10-> For example, I could walk for \apath{red}{$27s$} to this intersection.%
%
\item<11-> Or for \apath{red}{$83s$} to that one.%
%
\item<12-> Or for \apath{red}{$195s$} to that one.%
}%
%
\only<-23>{%
\item<13-> We have 3~choices: (a)~$\apath{red}{27s}\only<17->{+\apath{green}{213s}}$, (b)~$\apath{red}{83s}\only<18->{+\apath{green}{221s}}$, (c)~$\apath{red}{195s}\only<19->{+\apath{green}{362s}}$.%
%
\item<14-> For each intersection, we can compute the airline distance~(\apath{green}{as the crow flies}) to the~{\color{goal-color}{食堂}}.%
%
\item<15-> The actual walking distance can never be shorter than that.%
%
\item<16-> Combining the two distances for each node gives us a good impression of how useful going there would be.%
%
\item<23-> (a)~it is\dots%
}%
%
\only<-64>{%
\item<24-> We continue doing this procedure until the first path to the goal is found.%
%
\item<56-> We found a first complete path from {\color{start-color}{start}} to {\color{goal-color}{goal}}.%
%
\item<57-> It has the total length~\apath{red}{$268s$}.%
%
\item<58-> Now we need to only check the remaining paths that might be shorter, those that are at least as long can be ignored.%
%
\item<64-> We confirmed the shortest possible path and it takes~$268s$.%
}%
%
\item<65-> In the worst case, the \pgls{algoAstar} needs to \inQuotes{look} at all intersections~(nodes) and streets~(edges) once.%
%
\item<66-> Usually, it is quite efficient to find shortest paths on maps.%
%
\item<67-> If we look for shortest paths that do not visit any node twice, where the edge distances are not negative, where the number of choices per node is limited, and where the graph and path are not too big/long -- then this algorithm is very efficient.%
%
\end{itemize}%
}}{0.51}{0.1}%075}%
%
\locateGraphic{}{width=0.545\paperwidth}{\sharedDir/graphics/hfuu_campus_map_astar/hfuu_campus_map_astar}{-0.02}{0.15}%
%
\locateGraphic{2-6}{width=0.545\paperwidth}{\sharedDir/graphics/hfuu_campus_map_astar/hfuu_campus_map_astar_01_start_goal}{-0.02}{0.15}%
\locateGraphic{8-}{width=0.55045\paperwidth}{\sharedDir/graphics/hfuu_campus_map_astar/hfuu_campus_map_astar_03_overlay}{-0.02421}{0.1449}%
%
\locateGraphic{7-9}{width=0.545\paperwidth}{\sharedDir/graphics/hfuu_campus_map_astar/hfuu_campus_map_astar_02_paths}{-0.02}{0.15}%
%
\locateGraphic{10}{width=0.545\paperwidth}{\sharedDir/graphics/hfuu_campus_map_astar/hfuu_campus_map_astar_04a}{-0.02}{0.15}%
\locateGraphic{11}{width=0.545\paperwidth}{\sharedDir/graphics/hfuu_campus_map_astar/hfuu_campus_map_astar_04b}{-0.02}{0.15}%
\locateGraphic{12-16}{width=0.545\paperwidth}{\sharedDir/graphics/hfuu_campus_map_astar/hfuu_campus_map_astar_04c}{-0.02}{0.15}%
\locateGraphic{17}{width=0.545\paperwidth}{\sharedDir/graphics/hfuu_campus_map_astar/hfuu_campus_map_astar_04d}{-0.02}{0.15}%
\locateGraphic{18}{width=0.545\paperwidth}{\sharedDir/graphics/hfuu_campus_map_astar/hfuu_campus_map_astar_04e}{-0.02}{0.15}%
\locateGraphic{19}{width=0.545\paperwidth}{\sharedDir/graphics/hfuu_campus_map_astar/hfuu_campus_map_astar_04f}{-0.02}{0.15}%
%
\locateGraphic{20}{width=0.545\paperwidth}{\sharedDir/graphics/hfuu_campus_map_astar/hfuu_campus_map_astar_04g}{-0.02}{0.15}%
\locateGraphic{21}{width=0.545\paperwidth}{\sharedDir/graphics/hfuu_campus_map_astar/hfuu_campus_map_astar_04h}{-0.02}{0.15}%
\locateGraphic{22}{width=0.545\paperwidth}{\sharedDir/graphics/hfuu_campus_map_astar/hfuu_campus_map_astar_04i}{-0.02}{0.15}%
%
\locateGraphic{23}{width=0.545\paperwidth}{\sharedDir/graphics/hfuu_campus_map_astar/hfuu_campus_map_astar_05}{-0.02}{0.15}%
%
\locateGraphic{24}{width=0.545\paperwidth}{\sharedDir/graphics/hfuu_campus_map_astar/hfuu_campus_map_astar_06a}{-0.02}{0.15}%
\locateGraphic{25}{width=0.545\paperwidth}{\sharedDir/graphics/hfuu_campus_map_astar/hfuu_campus_map_astar_06b}{-0.02}{0.15}%
%
\locateGraphic{26}{width=0.545\paperwidth}{\sharedDir/graphics/hfuu_campus_map_astar/hfuu_campus_map_astar_06c}{-0.02}{0.15}%
\locateGraphic{27}{width=0.545\paperwidth}{\sharedDir/graphics/hfuu_campus_map_astar/hfuu_campus_map_astar_06d}{-0.02}{0.15}%
%
\locateGraphic{28}{width=0.545\paperwidth}{\sharedDir/graphics/hfuu_campus_map_astar/hfuu_campus_map_astar_07a}{-0.02}{0.15}%
\locateGraphic{29}{width=0.545\paperwidth}{\sharedDir/graphics/hfuu_campus_map_astar/hfuu_campus_map_astar_07b}{-0.02}{0.15}%
%
\locateGraphic{30}{width=0.545\paperwidth}{\sharedDir/graphics/hfuu_campus_map_astar/hfuu_campus_map_astar_07c}{-0.02}{0.15}%
%
\locateGraphic{31}{width=0.545\paperwidth}{\sharedDir/graphics/hfuu_campus_map_astar/hfuu_campus_map_astar_08a}{-0.02}{0.15}%
\locateGraphic{32}{width=0.545\paperwidth}{\sharedDir/graphics/hfuu_campus_map_astar/hfuu_campus_map_astar_08b}{-0.02}{0.15}%
\locateGraphic{33}{width=0.545\paperwidth}{\sharedDir/graphics/hfuu_campus_map_astar/hfuu_campus_map_astar_09a}{-0.02}{0.15}%
\locateGraphic{34}{width=0.545\paperwidth}{\sharedDir/graphics/hfuu_campus_map_astar/hfuu_campus_map_astar_09b}{-0.02}{0.15}%
%
\locateGraphic{35}{width=0.545\paperwidth}{\sharedDir/graphics/hfuu_campus_map_astar/hfuu_campus_map_astar_10a}{-0.02}{0.15}%
\locateGraphic{36}{width=0.545\paperwidth}{\sharedDir/graphics/hfuu_campus_map_astar/hfuu_campus_map_astar_10b}{-0.02}{0.15}%
\locateGraphic{37}{width=0.545\paperwidth}{\sharedDir/graphics/hfuu_campus_map_astar/hfuu_campus_map_astar_10c}{-0.02}{0.15}%
%
\locateGraphic{38}{width=0.545\paperwidth}{\sharedDir/graphics/hfuu_campus_map_astar/hfuu_campus_map_astar_10d}{-0.02}{0.15}%
\locateGraphic{39}{width=0.545\paperwidth}{\sharedDir/graphics/hfuu_campus_map_astar/hfuu_campus_map_astar_11a}{-0.02}{0.15}%
%
\locateGraphic{40}{width=0.545\paperwidth}{\sharedDir/graphics/hfuu_campus_map_astar/hfuu_campus_map_astar_11b}{-0.02}{0.15}%
%
\locateGraphic{41}{width=0.545\paperwidth}{\sharedDir/graphics/hfuu_campus_map_astar/hfuu_campus_map_astar_11c}{-0.02}{0.15}%
%
\locateGraphic{42}{width=0.545\paperwidth}{\sharedDir/graphics/hfuu_campus_map_astar/hfuu_campus_map_astar_12a}{-0.02}{0.15}%
\locateGraphic{43}{width=0.545\paperwidth}{\sharedDir/graphics/hfuu_campus_map_astar/hfuu_campus_map_astar_12b}{-0.02}{0.15}%
\locateGraphic{44}{width=0.545\paperwidth}{\sharedDir/graphics/hfuu_campus_map_astar/hfuu_campus_map_astar_12c}{-0.02}{0.15}%
\locateGraphic{45}{width=0.545\paperwidth}{\sharedDir/graphics/hfuu_campus_map_astar/hfuu_campus_map_astar_12d}{-0.02}{0.15}%
%
\locateGraphic{46}{width=0.545\paperwidth}{\sharedDir/graphics/hfuu_campus_map_astar/hfuu_campus_map_astar_13a}{-0.02}{0.15}%
\locateGraphic{47}{width=0.545\paperwidth}{\sharedDir/graphics/hfuu_campus_map_astar/hfuu_campus_map_astar_13b}{-0.02}{0.15}%
%
\locateGraphic{48}{width=0.545\paperwidth}{\sharedDir/graphics/hfuu_campus_map_astar/hfuu_campus_map_astar_13c}{-0.02}{0.15}%
%
\locateGraphic{49}{width=0.545\paperwidth}{\sharedDir/graphics/hfuu_campus_map_astar/hfuu_campus_map_astar_14a}{-0.02}{0.15}%
\locateGraphic{50}{width=0.545\paperwidth}{\sharedDir/graphics/hfuu_campus_map_astar/hfuu_campus_map_astar_14b}{-0.02}{0.15}%
\locateGraphic{51}{width=0.545\paperwidth}{\sharedDir/graphics/hfuu_campus_map_astar/hfuu_campus_map_astar_14c}{-0.02}{0.15}%
\locateGraphic{52}{width=0.545\paperwidth}{\sharedDir/graphics/hfuu_campus_map_astar/hfuu_campus_map_astar_14d}{-0.02}{0.15}%
%
\locateGraphic{53}{width=0.545\paperwidth}{\sharedDir/graphics/hfuu_campus_map_astar/hfuu_campus_map_astar_15a}{-0.02}{0.15}%
\locateGraphic{54}{width=0.545\paperwidth}{\sharedDir/graphics/hfuu_campus_map_astar/hfuu_campus_map_astar_15b}{-0.02}{0.15}%
\locateGraphic{55}{width=0.545\paperwidth}{\sharedDir/graphics/hfuu_campus_map_astar/hfuu_campus_map_astar_15c}{-0.02}{0.15}%
%
\locateGraphic{56}{width=0.545\paperwidth}{\sharedDir/graphics/hfuu_campus_map_astar/hfuu_campus_map_astar_16}{-0.02}{0.15}%
%
\locateGraphic{57}{width=0.545\paperwidth}{\sharedDir/graphics/hfuu_campus_map_astar/hfuu_campus_map_astar_17}{-0.02}{0.15}%
\locateGraphic{58}{width=0.545\paperwidth}{\sharedDir/graphics/hfuu_campus_map_astar/hfuu_campus_map_astar_18}{-0.02}{0.15}%
\locateGraphic{59}{width=0.545\paperwidth}{\sharedDir/graphics/hfuu_campus_map_astar/hfuu_campus_map_astar_19}{-0.02}{0.15}%
\locateGraphic{60}{width=0.545\paperwidth}{\sharedDir/graphics/hfuu_campus_map_astar/hfuu_campus_map_astar_20}{-0.02}{0.15}%
\locateGraphic{61}{width=0.545\paperwidth}{\sharedDir/graphics/hfuu_campus_map_astar/hfuu_campus_map_astar_21}{-0.02}{0.15}%
\locateGraphic{62}{width=0.545\paperwidth}{\sharedDir/graphics/hfuu_campus_map_astar/hfuu_campus_map_astar_22}{-0.02}{0.15}%
\locateGraphic{63-}{width=0.545\paperwidth}{\sharedDir/graphics/hfuu_campus_map_astar/hfuu_campus_map_astar_23}{-0.02}{0.15}%
%
\end{frame}%
%
\begin{frame}%
\frametitle{Transition to Hard Problems}%
\begin{itemize}%
\item So we first looked at an optimization problem that can easily be solved to optimality using high school mathematics.%
%
\item<2-> Then we looked at a problem that can be solved to optimality using an algorithm that you should have learned during your Bachelor's studies.%
%
\item<3-> What about \emph{hard} problems?%
%
\item<4-> What about problems where traditional Maths and algorithms do not really help us much?%
%
\item<5-> Here we go.%
\end{itemize}%
\end{frame}%
%
%
%
\begin{frame}%
\frametitle{Hard Problems~(1)}%
\begin{itemize}%
\item There exists a group of problems that are \emph{hard}.%
%
\item<2-> Actually, most of the problems I mention at the beginning of this talk belong to this group of hard problems.%
%
\item<3-> A hard problem cannot be solved both exactly and efficiently.%
%
\item<4-> Let's look at two quick examples.%
\end{itemize}%
\end{frame}%
%
\begin{frame}[t]%
\frametitle{Bin Packing}%
\begin{itemize}%
\item My car can carry $T$~kg of weight.%
%
\item<2-> I have $n$~objects, each with weight~$w_i$ for~$i$ in~$1$ to~$n$.%
%
\item<3-> How can I pack my car so that I can carry them from~$A$ to $B$ with the fewest possible hauls?\cite{ZLWvdBTW2024RLSOT2RBPPWIR,ZWvdBTLTW2024RLSFTDBPAANRFFFA}%
\end{itemize}%
%
\locateGraphic{}{width=0.8\paperwidth}{\sharedDir/graphics/bin_packing/bin_packing}{0.1}{0.47}%
\end{frame}%
%
\begin{frame}[t]%
\frametitle{Traveling Salesperson Problem}%
\begin{itemize}%
\item Find the shortest path that visits~$n$ locations and returns back to its origin.\cite{ABCC2006TTSPACS,LLRKS1985TTSPAGTOCO}%
\end{itemize}%
%
\locateGraphic{}{width=0.5\paperwidth}{\sharedDir/graphics/tsp_china/tsp_china}{0.25}{0.24}%
\end{frame}%
%
%
\def\makebigmarkformat{%
\def\bigmarkformat##1{%
{\alert{\textbf{\Large{##1}}}}%
}}%
\protected\gdef\bigmark#1#2{%
\def\bigmarkformat##1{##1}%
\only<#1>{\makebigmarkformat}%
\bigmarkformat{#2}%
}%
%
\begin{frame}[t]%
\frametitle{Hard Problems~(2)}%
\only<-3,20->{%
\begin{itemize}%
\item These are just two examples from a huge family of problems called~\npHard.%
%
\item<2-> What does that mean?%
\end{itemize}%
}%
%
%
\locate{3-}{\parbox{0.9\paperwidth}{\noindent%
\makeatletter%
\begin{mdframed}[style=@bestPracticeStyle]\Large{%
If an algorithm guarantees to always find the \bigmark{5-6}{best-possible} solution for an \npHard\ problem, then its \bigmark{22}{worst-case} runtime will \bigmark{9-19}{grow exponential} with the \bigmark{7-8}{input size~\only<-8,20->{\emph{s}}}\only<9-19>{\bigmark{9-19}{\emph{s}}}.%
}\end{mdframed}%
\makeatother%
}}{0.05}{0.7}%
%
\locateGraphic{4-9}{width=0.4\paperwidth}{\sharedDir/graphics/tsp_china/tsp_china}{0.06666666666}{0.05}%
\locateGraphic{4-9}{width=0.4\paperwidth}{\sharedDir/graphics/bin_packing/bin_packing}{0.53333333333}{0.25}%
%
\locate{6}{\parbox{0.4\paperwidth}{\noindent\centering\Large{%
\alert{shortest round-trip tour}%
}}}{0.06666666666}{0.65}%
\locate{6}{\parbox{0.4\paperwidth}{\noindent\centering\Large{%
\alert{packing with fewest hauls}%
}}}{0.53333333333}{0.65}%
%
\locate{8}{\parbox{0.4\paperwidth}{\noindent\centering\Large{%
\alert{number of cities to visit}%
}}}{0.06666666666}{0.65}%
\locate{8}{\parbox{0.4\paperwidth}{\noindent\centering\Large{%
\alert{number of items to pack}%
}}}{0.53333333333}{0.65}%
%
\locate{21}{\parbox{0.7\paperwidth}{\noindent\centering\textbf{\Large{\alert{%
If we want to get the optimal solution,\\it will take too long.%
}}}}}{0.15}{0.45}%
%
\locate{22}{\parbox{0.7\paperwidth}{\noindent\centering\textbf{\Large{%
If we want to get the optimal solution,\\it will \alert{sometimes/often/always} take too long.%
}}}}{0.15}{0.45}%
%
\locateGraphic{10}{width=0.8\paperwidth}{\sharedDir/graphics/function_growth/function_growth_01}{0.1}{0.05}%
\locateGraphic{11}{width=0.8\paperwidth}{\sharedDir/graphics/function_growth/function_growth_02}{0.1}{0.05}%
\locateGraphic{12}{width=0.8\paperwidth}{\sharedDir/graphics/function_growth/function_growth_03}{0.1}{0.05}%
\locateGraphic{13}{width=0.8\paperwidth}{\sharedDir/graphics/function_growth/function_growth_04}{0.1}{0.05}%
\locateGraphic{14}{width=0.8\paperwidth}{\sharedDir/graphics/function_growth/function_growth_05}{0.1}{0.05}%
\locateGraphic{15}{width=0.8\paperwidth}{\sharedDir/graphics/function_growth/function_growth_06}{0.1}{0.05}%
\locateGraphic{16}{width=0.8\paperwidth}{\sharedDir/graphics/function_growth/function_growth_07}{0.1}{0.05}%
\locateGraphic{17}{width=0.8\paperwidth}{\sharedDir/graphics/function_growth/function_growth_08}{0.1}{0.05}%
\locateGraphic{18}{width=0.8\paperwidth}{\sharedDir/graphics/function_growth/function_growth_09}{0.1}{0.05}%
\locateGraphic{19}{width=0.8\paperwidth}{\sharedDir/graphics/function_growth/function_growth_10}{0.1}{0.05}%
%
\end{frame}%
%
\global\let\sharedSlideTitle\relax%
%
%
%
\section{Metaheuristics}%
%
%
\begin{frame}%
\sharedframetitle{Metaheuristics}%
\parbox{0.42\paperwidth}{%
\begin{itemize}%
%
\item Because of this hardness of optimization, metaheuristic algorithms that follow the Trial-and-Error idea of iterative improvement have emerged.%
\item<2-> They drop the guarantee to find the optimal solution.%
\item<3-> They try to find good solutions within a feasible runtime.%
\item<4-> They (usually) start with random solutions.%
\item<5-> And then roughly follow this cycle.%
%
\end{itemize}%
}%
%
\locateGraphic{4}{width=0.5\paperwidth}{\sharedDir/graphics/metaheuristic_cycle/metaheuristic_cycle_01}{0.47}{0.3}%
%
\locateGraphic{5}{width=0.5\paperwidth}{\sharedDir/graphics/metaheuristic_cycle/metaheuristic_cycle_02}{0.47}{0.3}%
%
\locateGraphic{6}{width=0.5\paperwidth}{\sharedDir/graphics/metaheuristic_cycle/metaheuristic_cycle_03}{0.47}{0.3}%
%
\locateGraphic{7}{width=0.5\paperwidth}{\sharedDir/graphics/metaheuristic_cycle/metaheuristic_cycle_04}{0.47}{0.3}%
%
\locateGraphic{8}{width=0.5\paperwidth}{\sharedDir/graphics/metaheuristic_cycle/metaheuristic_cycle_05}{0.47}{0.3}%
%
\locateGraphic{9}{width=0.5\paperwidth}{\sharedDir/graphics/metaheuristic_cycle/metaheuristic_cycle_06}{0.47}{0.3}%
%
\locateGraphic{10}{width=0.5\paperwidth}{\sharedDir/graphics/metaheuristic_cycle/metaheuristic_cycle_07}{0.47}{0.3}%
%
\end{frame}%
%
\global\let\sharedSlideTitle\relax%
%
%
%
\gdef\sharedSlideTitle{Example:~RLS and the $(1+1)$~EA}%
%
%% needs command algobox, as in ffa slides
%
\begin{frame}[t]%
\sharedframetitle{RLS and the $(1+1)$~EA}%
\begin{itemize}%
\item Local search with $|S_i|=|N_i|=1$ is the simplest realization of the metaheuristic idea.%
\item<2-> \GlsFull{algoRLS} and the \glsFull{algoEAopo} work according the same pattern~(and differ only in their unary search operator~$\mathop{move}$)\cite{DJW2002OTAOTOPOEA,CPD2018TAMPARAOEA}.%
\item<8-> They accept the new solution if it is better or equally good compared to the current solution.%
\end{itemize}%
%
%
\locate{2}{%
\algobox{0.75}{\includegraphics[width=0.75\paperwidth]{\sharedDir/graphics/metaheuristic_algorithms/opoea/opoea_1}}%
}{0.125}{0.445}%
%
\locate{3}{%
\algobox{0.75}{\includegraphics[width=0.75\paperwidth]{\sharedDir/graphics/metaheuristic_algorithms/opoea/opoea_2}}%
}{0.125}{0.445}%
%
\locate{4}{%
\algobox{0.75}{\includegraphics[width=0.75\paperwidth]{\sharedDir/graphics/metaheuristic_algorithms/opoea/opoea_3}}%
}{0.125}{0.445}%
%
\locate{5}{%
\algobox{0.75}{\includegraphics[width=0.75\paperwidth]{\sharedDir/graphics/metaheuristic_algorithms/opoea/opoea_4}}%
}{0.125}{0.445}%
%
\locate{6}{%
\algobox{0.75}{\includegraphics[width=0.75\paperwidth]{\sharedDir/graphics/metaheuristic_algorithms/opoea/opoea_5}}%
}{0.125}{0.445}%
%
\locate{7}{%
\algobox{0.75}{\includegraphics[width=0.75\paperwidth]{\sharedDir/graphics/metaheuristic_algorithms/opoea/opoea_6}}%
}{0.125}{0.445}%
%
\locate{8}{%
\algobox{0.75}{\includegraphics[width=0.75\paperwidth]{\sharedDir/graphics/metaheuristic_algorithms/opoea/opoea_7}}%
}{0.125}{0.445}%
%
\locate{9}{%
\algobox{0.75}{\includegraphics[width=0.75\paperwidth]{\sharedDir/graphics/metaheuristic_algorithms/opoea/opoea_8}}%
}{0.125}{0.445}%
%
\locate{10}{%
\algobox{0.75}{\includegraphics[width=0.75\paperwidth]{\sharedDir/graphics/metaheuristic_algorithms/opoea/opoea_9}}%
}{0.125}{0.445}%
%
\locate{11-}{%
\algobox{0.75}{\includegraphics[width=0.75\paperwidth]{\sharedDir/graphics/metaheuristic_algorithms/opoea/opoea}}%
}{0.125}{0.445}%
%
\end{frame}%
%
\global\let\sharedSlideTitle\relax%
%
%
%
\xdef\sharedSlideTitle{Example:~Simulated Annealing}%
%
%% needs command algobox, as in ffa slides
%
\begin{frame}%
\sharedframetitle{Simulated Annealing}%
\parbox{0.41\paperwidth}{%
\begin{itemize}%
\item \glsFull{algoSA}\cite{KGV1983OBSA} is a local search that accepts also worsening moves, but with a probability that decreases over time and with the difference in solution quality.%
\item<3-> ${\mathfrak{R}}_0^1$~be a random number uniformly distributed in~$[0,1)$.%
\item<6-> High temperatures~$T$ increase the acceptance probability~$P$ of worse solutions, low~$T$ reduce it.%
\item<7-> $P$~is regulated by temperature schedule with parameter\only<10->{s}~$T_0$\only<10->{ and~$\varepsilon$}.%
\end{itemize}%
}%
\uncover<11->{\begin{itemize}%
\item \gls{algoSA} also remembers best-so-far solution~$x_B$ and its objective value~$y_B$, because it could get lost.%
\end{itemize}%
}%
%
%
\locate{2}{%
\algobox{0.5}{\includegraphics[width=0.5\paperwidth]{\sharedDir/graphics/metaheuristic_algorithms/sa/sa_01}}%
}{0.44}{0.24}%
%
\locate{3}{%
\algobox{0.5}{\includegraphics[width=0.5\paperwidth]{\sharedDir/graphics/metaheuristic_algorithms/sa/sa_02}}%
}{0.44}{0.24}%
%
\locate{4}{%
\algobox{0.4}{\includegraphics[width=0.4\paperwidth]{\sharedDir/graphics/metaheuristic_algorithms/sa/sa_02}}%
}{0.44}{0.07}%
%
\locate{4}{%
\resizebox{0.5\paperwidth}{!}{\bgroup%
\fboxsep=2pt%
\fboxrule=2pt%
\fcolorbox{black}{white}{%
\includegraphics[width=0.5\paperwidth]{\sharedDir/graphics/metaheuristic_algorithms/sa/exp}%
}%
\egroup}%%
}{0.46}{0.41}%
%
\locate{5}{%
\algobox{0.5}{\includegraphics[width=0.5\paperwidth]{\sharedDir/graphics/metaheuristic_algorithms/sa/sa_03}}%
}{0.44}{0.24}%
%
\locate{6}{%
\algobox{0.5}{\includegraphics[width=0.5\paperwidth]{\sharedDir/graphics/metaheuristic_algorithms/sa/sa_04}}%
}{0.44}{0.24}%
%
\locate{7}{%
\algobox{0.5}{\includegraphics[width=0.5\paperwidth]{\sharedDir/graphics/metaheuristic_algorithms/sa/sa_05}}%
}{0.44}{0.24}%
%
\locate{8}{%
\algobox{0.5}{\includegraphics[width=0.5\paperwidth]{\sharedDir/graphics/metaheuristic_algorithms/sa/sa_06}}%
}{0.44}{0.24}%
%
\locate{9}{%
\algobox{0.5}{\includegraphics[width=0.5\paperwidth]{\sharedDir/graphics/metaheuristic_algorithms/sa/sa_07}}%
}{0.44}{0.24}%
%
\locate{10}{%
\algobox{0.5}{\includegraphics[width=0.5\paperwidth]{\sharedDir/graphics/metaheuristic_algorithms/sa/sa_08}}%
}{0.44}{0.24}%
%
\locate{11}{%
\algobox{0.5}{\includegraphics[width=0.5\paperwidth]{\sharedDir/graphics/metaheuristic_algorithms/sa/sa_09}}%
}{0.44}{0.24}%
%
\locate{12}{%
\algobox{0.5}{\includegraphics[width=0.5\paperwidth]{\sharedDir/graphics/metaheuristic_algorithms/sa/sa_10}}%
}{0.44}{0.24}%
%
\locate{13}{%
\algobox{0.5}{\includegraphics[width=0.5\paperwidth]{\sharedDir/graphics/metaheuristic_algorithms/sa/sa}}%
}{0.44}{0.24}%
%
\end{frame}%
%
\global\let\sharedSlideTitle\relax%
%
%
%
%
\begin{frame}[t]%
\sharedframetitle{Components needed for Metaheuristic Optimization}%
%
\begin{itemize}%
%
\item We have a space~\searchSpace\ in which we search for potential solutions~\sespel.%
%
\item<2-> \searchSpace~is often very/exponentially large in the input size.%
%
\item<3-> We repeatedly sample elements~\sespel\ from this space~\searchSpace, so we need operators to do that\only<-3>{.}\only<4->{:%
\begin{itemize}%
\item a \emph{nullary} operator usually returns a randomly chosen point from~$\searchSpace$.\uncover<5->{ %
This is often used to get starting points for the search.}%
%
\item<6-> a \emph{unary} operator receives an existing point as input and creates a slightly modified copy.\uncover<7->{ %
This is often used to explore the neighborhood of good solutions.}%
%
\item<8-> a \emph{binary} operator receives a two existing points as input and creates a new point with elements from both of these \inQuotes{parents}.\uncover<9->{ %
This is often used to combine traits from good solutions.}%
%
\end{itemize}%
}%
%
\item<10-> We evaluate the quality of solutions using the objective function~$\objFun:\searchSpace\mapsto\realNumbers$, which usually is subject to minimization~(smaller values are better).%
%
\item<11-> These are the usual main ingredients of \gls{metaheuristic} optimization.%
%
\end{itemize}%
%
\end{frame}%
%
\global\let\sharedSlideTitle\relax%
%
%
%
\begin{frame}[t]%
\frametitle{Types of Metaheuristics}%
%
\begin{itemize}%
%
\item We can build many interesting algorithms just using such components~(and sometimes a bit more problem-specific information).
%
\item<2-> Local Search\only<3->{:%
\begin{itemize}%
\item Keep one solution~$x_c$, sample at least one new solution~$x_n$ from its neighborhood in each step, move to better solutions as discovered.%
%
\only<-5>{%
\item<4-> {\color{green!60!black}Advantage}: Fast%
\item<5-> {\color{red!60!black}Disadvantage}: May get stuck at a local optimum%
}%
\end{itemize}%
}%
%
\item<6-> Population-based Methods\only<7->{:%
\begin{itemize}%
\item Maintain population~$S$ of $|S|=\mu\geq1$ solutions, sample set~$N$ of $|N|=\lambda\geq1$ new solutions in each step by creating modified copies of the elements in~$S$, then update $S$.%
\only<-9>{%
\item<8-> {\color{green!60!black}Advantage}: Less likely to get stuck at local optimum.%
\item<9-> {\color{red!60!black}Disadvantage}: Slow.%
}%
\end{itemize}%
}%
%
\item<10-> Model-based Methods\only<11->{:%
\begin{itemize}%
\item Maintain a stochastic model~$M$ (usually a multivariate probility distribution) of what good solutions could look like, sample set~$N$ of $|N|=\lambda\geq1$ new solutions based on~$M$ in each step, then update~$M$.%
\only<-13>{%
\item<12-> {\color{green!60!black}Advantage}: Sometimes really good performance.%
\item<13-> {\color{red!60!black}Disadvantage}: Complicated, sampling and model update may be slow.%
}%
\end{itemize}%
}%
%
\item<14-> Hybrid Algorithms\only<15->{:%
\begin{itemize}%
\item Usually combine model-based or population-based approach with local search: Each new solution from model/population is refined with local search.%
\only<-17>{%
\item<16-> {\color{green!60!black}Advantage}: Fast (because local search), good global search (because of population/model)%
\item<17-> {\color{red!60!black}Disadvantage}: usually few.%
}%
\end{itemize}%
}%
%
\end{itemize}%
%
\end{frame}%
%
\section{Fields where Metaheuristics can be Applied}%
%
\begin{frame}%
\frametitle{Types of Search Spaces / Decision Variables}%
%
\begin{itemize}%
\item Discrete problems have solutions that can be expressed as bit strings of length~$n$, i.e., search in~$\{0,1\}^n$.%
%
\item<2-> Combinatorial problems have solutions that can be expressed as permutations of first $n$~natural numbers.%
%
\item<3-> Integer problems have solutions that can be expressed as integer vectors~$\subset\integerNumbers^n$.%
%
\item<4-> Continuous problems have solutions that can be expressed as real vectors~$\subset\realNumbers^n$.%
%
\item<5-> Mixed problems: usually discrete + integer + continuous.%
%
\item<6-> Genetic Programming usually has solutions that can be expressed as tree data structures with different types of nodes and a maximum depth~(imagine mathematical expressions like~$(5x+3)/3$, which have inherent tree nature).
\end{itemize}%
\end{frame}%
%
\begin{frame}%
\frametitle{Mode of the Objective Function(s)}%
%
\begin{itemize}%
\item The objective function~$f(x)$ tells us how good a solution~$x$ is; the smaller~$f(x)$, the better.%
%
\item<2-> $f$~can be a mathematical function with properties such as differentiability, etc.%
%
\item<3-> $f$~can be discrete function, like on the TSP.%
%
\item<4-> $f$~may require us to run complex simulations, e.g., if $x$~is some engineering design or a policy that needs to be simulated.%
%
\item<5-> $f$~usually is static and does not change{\dots}%
%
\item<6-> {\dots}but sometimes, $f$~is dynamic and may change while we try to optimize it.%
%
\item<7-> Sometimes, we face a multi-objective problem, where we do not have a single objective function, but multiple functions, i.e., a vector function~$\vec{f}(x)$, and we try to trade-off different goals; for example speed vs.\ cost.%
\end{itemize}%
\end{frame}%
%
\endPresentation%
\end{document}%%
\endinput%
%
